\section{Labelled pieces of algebraic pushouts}\label{sec:pieces}

We fix $1\leq n\leq\infty$ throughout this section.  The arguments concern
actual subcategories and actual pushouts.  This precision will allow us to
apply the simplicial gluing lemma even when an attaching functor identifies
objects or cells.

The sieve and nerve constructions have established antecedents in
\cite[Propositions~5.3 and~5.5]{AraMaltsiniotis2014}. We give the
calculations explicitly to identify the labels, fibers and intersections
after an arbitrary attachment.

A full sub-$n$-category of $C$ on a set of objects $S$ has exactly those
objects and contains every positive-dimensional cell whose $0$-source
and $0$-target belong to $S$.  It is a \emph{sieve} if a
$1$-cell $x\to y$ with $y\in S$ forces $x\in S$, and a \emph{cosieve} if a
$1$-cell $x\to y$ with $x\in S$ forces $y\in S$.  A sieve is equivalently
the fiber over $0$ of a strict functor $C\to\Disk{1}$; a cosieve is
equivalently the fiber over $1$.  Indeed, the indicated object assignment
extends uniquely to a strict functor: each nonempty hom in the target is
terminal, and the sieve or cosieve condition excludes an arrow from a
$1$-labelled object to a $0$-labelled object.  These fibers are full, since
all cells between objects with the same label necessarily map to identities
at that label.

The following fiber adjoints are those of
\cite[Lemma~2.8]{AraMaltsiniotis2014}.

\begin{lemma}[Fibers and objects preserve colimits]\label{lem:fibers}
For $\epsilon\in\{0,1\}$, the functor
\[
 F_\epsilon:\nCat/\Disk{1}\longrightarrow\nCat,
 \qquad (C\xrightarrow{p}\Disk{1})\longmapsto C_\epsilon=p^{-1}(\epsilon),
\]
preserves all small colimits.  The object functor
$\operatorname{Ob}:\nCat\to\mathrm{Set}$ also preserves all small colimits.
Both assertions hold directly for $n=\infty$.
\end{lemma}

\begin{proof}
We construct a right adjoint $Q_0$ to $F_0$.  Given $B$, put all its
objects over $0$ and adjoin a new object $*$ over $1$.  Keep the hom
categories $B(b,b')$, and set
\[
 Q_0(B)(b,*)=\mathbf{1},\qquad
 Q_0(B)(*,b)=\varnothing,\qquad
 Q_0(B)(*,*)=\mathbf{1}.
\]
Here the homs are strict $(n-1)$-categories for finite $n$, sets when
$n=1$, and strict $\omega$-categories when $n=\infty$; $\mathbf{1}$ and
$\varnothing$ denote their terminal and initial objects.  Compositions
inside $B$ are retained.  Every remaining composition either has empty
domain or has terminal codomain, so is uniquely determined.  Associativity,
units, and interchange hold by the laws in $B$ or by uniqueness.  This
defines $Q_0(B)$ and its map to $\Disk{1}$ in all the stated dimensions.

For $p:C\to\Disk{1}$, a functor $C_0\to B$ extends uniquely over
$\Disk{1}$ to $C\to Q_0(B)$.  All objects of $C_1$ must go to $*$; their
internal cells and the cells from $C_0$ to $C_1$ have unique images in the
specified terminal homs.  There are no cells from $C_1$ to $C_0$, since
their $1$-boundaries would map to an arrow $1\to0$.  Restriction and
extension are inverse and natural, giving
\[
 \Hom_{\nCat/\Disk{1}}(C,Q_0(B))
 \cong \Hom_{\nCat}(C_0,B).
\]
Thus $F_0$ is a left adjoint.  For $F_1$, place $B$ over $1$ and adjoin
a new initial object over $0$, with terminal homs from that object into
$B$, empty reverse homs, and terminal endomorphism hom.  The same
restriction argument gives a right adjoint $Q_1$.  In particular these
constructions include $B=\varnothing$ and prove preservation of the empty
colimit as well.

For a set $S$, let $K(S)$ be the ordinary chaotic category on $S$, with
one arrow for each ordered pair of objects, regarded as an $n$-category
with identity cells above dimension one.  Every object function
$\operatorname{Ob}C\to S$ has a unique strict extension $C\to K(S)$.
Consequently $\operatorname{Ob}$ is left adjoint to $K$ and preserves
colimits.
\end{proof}

The next elementary property of orientals is needed before taking a
colimit of simplex pieces.  It is not enough merely to know their sets
of vertices.

\begin{lemma}[Full oriental segments]\label{lem:oriental-segments}
Let $I$ be an initial or a final segment of $[m]=\{0,\ldots,m\}$.
If $I$ is nonempty and has $q+1$ vertices, the face inclusion identifies
$\trunc\Orient{q}$ with the full sub-$n$-category of
$\trunc\Orient{m}$ on $I$.  If $I$ is empty, that full subcategory is
empty.  Here and below $\trunc=\id$ when $n=\infty$.
\end{lemma}

\begin{proof}
Recall the oriental presentation from Section~\ref{sec:foundations}:
its atoms are indexed by strictly increasing tuples, with $0$-source
the first vertex and $0$-target the last vertex.  All cells are finite
composites and identities of these atoms.  In particular every $1$-cell
has nondecreasing endpoints.

Suppose a cell has $0$-source $a$ and $0$-target $b$.  In any expression
for it, every occurring atom has all its vertices in $[a,b]$.  This
follows by induction on the expression.  For an atom it is the tuple
description.  In a $0$-composite the factors have endpoints $(a,c)$ and
$(c,b)$, with $a\leq c\leq b$.  In a composition in positive dimension
both factors have the same $0$-endpoints as the composite.  Identities
preserve these endpoints, and an identity on a vertex has support at
that vertex.  These observations prove the induction.  If $a,b\in I$,
the whole interval $[a,b]$ lies in $I$, so the expression uses only the
atoms of that face.  Thus the face generates every cell on those
objects.

For finite $n\geq1$, reflection leaves objects unchanged and every
cell of dimension at most $n$ is represented by a cell before
reflection.  The same argument applied to a representative proves
surjectivity onto the full subcategory after reflection; higher cells
are identities.  For injectivity, the inclusion $I\hookrightarrow[m]$
has a monotone retraction: clamp to the last vertex of $I$ for an
initial segment, or to its first vertex for a final segment.  By
cosimplicial functoriality it induces a retraction on orientals, and
then on their reflections.  The face map is therefore injective on
every cell set.  This proves the full-subcategory assertion.  If $I$
is empty there are no objects and hence no cells; no retraction is
needed.
\end{proof}

\begin{lemma}[Poset pieces]\label{lem:poset-pieces}
There is a canonical identification $\operatorname{Ob}\U{P}=P$ for
every poset $P$.  If $S\subset P$ is a sieve, the map
$\U{S}\to\U{P}$ is the inclusion of the full $0$-fiber of the membership
label $\U{P}\to\Disk{1}$.  If $S$ is a cosieve, it is the inclusion
of the full $1$-fiber.  These identifications are natural for maps of
labelled posets and include empty posets and empty fibers.
\end{lemma}

\begin{proof}
Use the canonical simplex presentation
\begin{equation}\label{eq:poset-simplex-colimit}
 \U{P}=\colim_{\sigma:\Delta[m]\to\OrdN P}\trunc\Orient{m}.
\end{equation}
By Lemma~\ref{lem:fibers}, its object set is the colimit of the vertex
sets $[m]$.  Sending the $j$th vertex of a simplex $\sigma$ to
$\sigma(j)$ identifies this colimit with $P$: each element occurs as
a zero-simplex, and the face inclusions from zero-simplices identify
every vertex with that occurrence.

For a sieve $S$, its characteristic map $P\to[1]$ assigns $0$ on $S$
and $1$ on its complement.  It is monotone and induces the displayed
label on $\U{P}$, since $\U{[1]}=\Disk{1}$.  The presentation
\eqref{eq:poset-simplex-colimit} is consequently a colimit in
$\nCat/\Disk{1}$ as well: slice colimits have the underlying colimits
in $\nCat$.  In its $\sigma$-piece the vertices mapping to $S$ form an
initial segment $I_\sigma$.  By Lemma~\ref{lem:oriental-segments},
the $0$-fiber of that piece is its truncated segment oriental.
Taking fibers commutes with the colimit by Lemma~\ref{lem:fibers}.

We verify that this colimit of segment orientals is canonically
$\U{S}$.  Restricting each nonempty simplex segment to $S$ gives a
map from its segment oriental to $\U{S}$; the empty segment has the
unique such map.  These maps are compatible with every simplex
operator.  In the reverse direction, each simplex entirely in $S$
supplies its whole oriental to the fiber colimit.  The two composites
are identities.  On a simplex of $S$ this is immediate.  On a
nonempty segment of an arbitrary simplex $\sigma$, it follows from
the face inclusion selecting $I_\sigma$, which is a morphism in the
simplex category of $\OrdN P$; on an empty segment there is nothing
to check.  The isomorphism so obtained is the canonical map
$\U{S}\to\U{P}$, followed into its fiber.  Fibers are full sieves,
which proves the claim.  The argument for a cosieve uses label $1$
on $S$, final vertex segments, and $F_1$ throughout.  All constructions
commute with maps preserving the labels.
\end{proof}

Here is the form of the fiber calculation that we will use in pushouts.
It also records precisely which hypotheses are needed of a label on an
arbitrary attaching target.

\begin{proposition}[Labelled attachments]\label{prop:labelled-pushout}
Let $A\subset P$ be a full subposet, let $a:\U{A}\to X$ be any strict
functor, and form
\[
 Y=X\amalg_{\U{A}}\U{P}.
\]
Suppose a monotone map $\lambda:P\to[1]$ and a strict functor
$\xi:X\to\Disk{1}$ satisfy
\[
 \xi a=U(\lambda|_A):\U{A}\longrightarrow\Disk{1},
 \qquad \U{[1]}=\Disk{1}.
\]
Write
$P_\epsilon=\lambda^{-1}(\epsilon)$,
$A_\epsilon=A\cap P_\epsilon$, and
$X_\epsilon=\xi^{-1}(\epsilon)$.  Then the induced label
$Y\to\Disk{1}$ has actual full fibers
\begin{equation}\label{eq:pushout-fiber}
 Y_\epsilon\cong
 X_\epsilon\amalg_{\U{A_\epsilon}}\U{P_\epsilon}.
\end{equation}
Moreover, independently of these labels,
\begin{equation}\label{eq:pushout-objects}
 \operatorname{Ob}Y\cong
 \operatorname{Ob}X\amalg_A P
 \cong\operatorname{Ob}X\amalg(P\setminus A).
\end{equation}
For two such pairs of compatible labels, write
$\bar\lambda_1,\bar\lambda_2:Y\to\Disk{1}$ for the induced labels.
Their $0$-fibers intersect in the $0$-fiber of
$\max\circ(\bar\lambda_1,\bar\lambda_2)$; their $1$-fibers intersect
in the $1$-fiber of $\min\circ(\bar\lambda_1,\bar\lambda_2)$.
Formula~\eqref{eq:pushout-fiber}, using the corresponding maximum or
minimum on both original legs, computes these intersections as well.
\end{proposition}

\begin{proof}
The compatibility assumption makes the defining pushout a pushout in
$\nCat/\Disk{1}$.  Applying $F_\epsilon$ and then
Lemmas~\ref{lem:fibers} and~\ref{lem:poset-pieces} gives
\eqref{eq:pushout-fiber}.  This identifies the canonical map from
the displayed pushout to $Y$ with the inclusion of a full fiber.
The object formula follows from the object part of Lemma~\ref{lem:fibers}
and Lemma~\ref{lem:poset-pieces}.  Its last isomorphism uses only the
injectivity of $A\hookrightarrow P$: old vertices of $P$ are replaced
by their images in $X$, while vertices in $P\setminus A$ remain distinct
new objects.  There is no injectivity requirement on $\operatorname{Ob}a$.

Two labels on $Y$ give a functor
$Y\to\Disk{1}\times\Disk{1}$.  The monotone functions
$\max,\min:[1]\times[1]\to[1]$ are strict functors, so their
composites are labels on $Y$.  An object has maximum label $0$ exactly
when both labels are $0$, and minimum label $1$ exactly when both
labels are $1$.  Full subcategories on the same objects are equal,
which proves the assertions about intersections.  The composite
labels restrict to the corresponding maximum or minimum labels on
both legs of the original pushout, so the same fiber formula applies.
\end{proof}

For use in the next two sections, we spell out both cover consequences.

\begin{corollary}[Sieve and cosieve covers after attachment]
\label{cor:attachment-covers}
With the notation of Proposition~\ref{prop:labelled-pushout}, the
following are canonical identifications of subcategories of $Y$.
\begin{enumerate}[label=\textup{(\roman*)}]
\item Suppose $P_1,P_2$ are sieves of $P$, contain $A$, and cover $P$.
Then the subcategories
\[
 Y_i=X\amalg_{\U{A}}\U{P_i}\quad(i=1,2)
\]
are sieves covering $Y$ on objects, and
\[
 Y_1\cap Y_2
 =X\amalg_{\U{A}}\U{P_1\cap P_2}.
\]
\item Suppose $T_1,T_2$ are cosieves of $P$, cover $P$, satisfy
$A\subset T_1$, and satisfy $A\cap T_2=\varnothing$.  Then
\[
 Y_1=X\amalg_{\U{A}}\U{T_1},\qquad Y_2=\U{T_2}
\]
are cosieves covering $Y$ on objects, and
$Y_1\cap Y_2=\U{T_1\cap T_2}$.
\end{enumerate}
If $A$ itself is a sieve of $P$, the canonical map $X\to Y$ is the
inclusion of a full sieve.
\end{corollary}

\begin{proof}
For (i), label membership in each $P_i$ by $0$ and label every object
of $X$ by $0$.  Apply Proposition~\ref{prop:labelled-pushout} to the
two labels and their maximum.  For (ii), label membership in each
$T_i$ by $1$; use constant label $1$ on $X$ for $T_1$ and constant
label $0$ on $X$ for $T_2$.  These agree on the attaching source by
the stated conditions.  The $1$-fiber of the second old-side label
is empty, and so is that of the minimum of the two old-side labels.
The same is true on their attaching sources.  Formula
\eqref{eq:pushout-fiber} therefore gives $\U{T_2}$ and
$\U{T_1\cap T_2}$ exactly.  In both cases
\eqref{eq:pushout-objects} proves the object cover.  Finally, when
$A$ is a sieve, label $A$ and $X$ by $0$ and $P\setminus A$ by $1$.
The $0$-fiber formula is $X\amalg_{\U{A}}\U{A}=X$.
\end{proof}

In particular, an attaching functor may identify two old objects,
collapse an old arrow to an identity, or identify parallel old cells.
These operations cause no ambiguity in the cover calculations: the
labels agree on every old cell before the pushout is taken.  The
fiber adjunctions compute the resulting subcategories after all those
identifications.

\begin{lemma}[Street nerves of two-piece covers]\label{lem:nerve-cover}
Let $C_1,C_2\subset C$ be two full sieves which cover the objects of
$C$, or two full cosieves which cover its objects.  Put
$C_{12}=C_1\cap C_2$.  The square
\[
\begin{tikzcd}
 \Street C_{12} \ar[r,hook] \ar[d,hook] & \Street C_2 \ar[d,hook] \\
 \Street C_1 \ar[r,hook] & \Street C
\end{tikzcd}
\]
is an ordinary pushout of simplicial sets.  All four maps are
monomorphisms, so it is also a homotopy pushout in the Kan--Quillen
model structure.
\end{lemma}

\begin{proof}
A full inclusion induces a monomorphism of Street nerves, and the
intersection of the two nerves is $\Street C_{12}$: a functor whose
image belongs to both subcategories has a unique factorization through
their full intersection.  It remains to prove that every simplex
$z:\Orient{m}\to C$ belongs to one of the two nerves.  In the sieve
case choose a piece containing $z(m)$.  For each vertex $j\leq m$,
the oriental has an arrow $j\to m$, whose image shows that $z(j)$
belongs to that piece.  Fullness then puts every cell of the simplex
in the same piece.  In the cosieve case choose a piece containing
$z(0)$ and use the arrows $0\to j$.  This includes $m=0$ and proves
the union statement degree by degree.  A union of simplicial subsets
is their ordinary pushout over the intersection.  Since its legs
are monomorphisms, the homotopy-pushout assertion follows from
Lemma~\ref{lem:simplicial-gluing}.
\end{proof}
